% speaker_notes.tex -- separate speech notes for the Beamer defence deck.
% Compile separately from main.tex to produce speaker_notes.pdf.

\documentclass[11pt,a4paper]{article}

\usepackage[T1]{fontenc}
\usepackage{lmodern}
\usepackage[margin=22mm]{geometry}
\usepackage{xcolor}
\usepackage{enumitem}
\usepackage{amsmath}
\usepackage{amssymb}
\usepackage[hidelinks]{hyperref}

\definecolor{AcademicNavy}{HTML}{17324D}
\definecolor{AcademicMuted}{HTML}{5F6B76}
\definecolor{AcademicRule}{HTML}{CBD5E1}

\setlength{\parindent}{0pt}
\setlength{\parskip}{4pt}
\emergencystretch=2em
\setlist[itemize]{leftmargin=*,topsep=2pt,itemsep=2pt}
\setlist[enumerate]{leftmargin=*,topsep=2pt,itemsep=2pt}

\newcounter{slideno}
\newenvironment{slidenote}[1]{%
  \refstepcounter{slideno}%
  \section*{\color{AcademicNavy}Slide \theslideno. #1}%
  \addcontentsline{toc}{section}{Slide \theslideno. #1}%
}{%
  \vspace{0.6em}%
  {\color{AcademicRule}\hrule}%
  \vspace{0.8em}%
}

\newcommand{\Purpose}{\vspace{0.2em}\textbf{Purpose:}\par}
\newcommand{\WhatToSay}{\vspace{0.2em}\textbf{What to say:}\par}
\newcommand{\KeyPoints}{\vspace{0.2em}\textbf{Key points:}\par}
\newcommand{\Transition}{\vspace{0.2em}\textbf{Transition to next slide:}\par}

\title{\color{AcademicNavy}Speaker Notes\\
AI-Guided Prediction of Thermoelastic Wave Propagation in Geological Media}
\author{Diploma defence speech notes}
\date{Bachelor's diploma defence, 2026}

\begin{document}

\maketitle
\tableofcontents
\newpage

\begin{slidenote}{AI-Guided Prediction of Thermoelastic Wave Propagation in Geological Media}
\Purpose
Open the defence, introduce the team, and set the research-prototype scope.

\WhatToSay
Good morning. The topic of our diploma work is \textit{AI-Guided Prediction of Thermoelastic Wave Propagation in Geological Media}. The work was prepared by Nurshanov Dias, Askarov Insar, Kassymkhanov Timur, and Kaidrakhmetova Dinara under the supervision of Dr. Bakhyt N. Alipova. The main idea is to study how artificial intelligence can support directional prediction of thermoelastic wave behaviour in simplified geological media. From the start, I would like to define the scope clearly: this is a research prototype and a comparative framework. It is not a field-calibrated geophysical simulator and it is not proposed as a replacement for full numerical modelling.

\KeyPoints
\begin{itemize}
  \item The defence concerns AI-assisted prediction of thermoelastic wave propagation.
  \item The task is source--probe directional prediction in simplified geological media.
  \item The result is a comparative research framework, not an industrial simulator.
\end{itemize}

\Transition
I will begin with the motivation and aim of the work.
\end{slidenote}

\begin{slidenote}{Motivation and Aim}
\Purpose
Explain why the topic matters and state the aim without overclaiming the prototype.

\WhatToSay
Temperature changes in rocks can affect stress, deformation, and wave propagation. This is relevant for geothermal systems, underground engineering, and geophysical interpretation, where thermal and mechanical effects can interact. The aim of the thesis is to develop and justify an AI-guided framework for directional prediction of thermoelastic wave propagation in geological media under simplified effective-medium assumptions. The important word here is ``framework''. We are not trying to model a real field site with all geological complexity. Instead, we create a controlled environment where different AI surrogate models can be compared under the same source--probe prediction task.

\KeyPoints
\begin{itemize}
  \item Thermal changes can create mechanical effects in rock.
  \item The aim is a directional prediction framework under simplified assumptions.
  \item The scope is a research prototype, not a field-validated simulator.
\end{itemize}

\Transition
Next, I will show the geological media and effective parameters used in the model.
\end{slidenote}

\begin{slidenote}{Geological Media We Model}
\Purpose
Show how geological materials enter the prototype through effective scalar parameters.

\WhatToSay
This slide shows representative effective parameters for four geological media: granite, basalt, sandstone, and limestone. The table includes density, compressional wave velocity, thermal conductivity, heat capacity, and thermal expansion where available. These values are used as effective scalar parameters under a locally homogeneous isotropic approximation. This means the model does not explicitly resolve pores, cracks, bedding, fluids, or anisotropy. The diagram on the right shows the simplified two-dimensional plane-strain domain with a thermal source point and a probe point. The line between them defines the direction of prediction.

\KeyPoints
\begin{itemize}
  \item Rock types differ in mechanical and thermal properties.
  \item The prototype uses effective isotropic material parameters.
  \item The source--probe geometry defines the directional query.
\end{itemize}

\Transition
After defining the materials, I will introduce the coupled thermoelastic system.
\end{slidenote}

\begin{slidenote}{Coupled Thermoelastic System}
\Purpose
Present the theoretical backbone of the work in a compact, defensible way.

\WhatToSay
The coupled thermoelastic system combines three ideas. First, the constitutive law relates stress to elastic strain and temperature perturbation. The term $\gamma=3K\alpha$ represents thermoelastic coupling, where $K$ is the bulk modulus and $\alpha$ is thermal expansion. Second, the equation of motion links stress gradients to acceleration and displacement. Third, the heat equation describes the evolution of temperature and includes coupling with volumetric deformation. In physical terms, heat changes temperature, temperature produces thermal strain, strain contributes to stress, and stress drives displacement and wave-like response.

\KeyPoints
\begin{itemize}
  \item Stress contains elastic and thermal contributions.
  \item Motion is driven by stress gradients.
  \item Heat transfer and deformation are coupled through the thermoelastic term.
\end{itemize}

\Transition
Now I will translate this physical model into the two-dimensional directional prediction task.
\end{slidenote}

\begin{slidenote}{2D Directional Prediction Task}
\Purpose
Define the practical prediction query and make the mathematical scope explicit.

\WhatToSay
The practical task is two-dimensional directional prediction. The model domain is $\Omega\subset\mathbb{R}^{2}$, and we assume plane strain, small deformation, and linear isotropic thermoelasticity. The primary fields are temperature $T(x,t)$ and in-plane displacement $u(x,t)=(u,v)$. A source point $x_s$ and a probe point $x_p$ define the vector $d=x_p-x_s$, the distance $L$, the unit direction $\hat d$, and the azimuth angle $\phi$. The model mapping $\mathcal{F}:X\to\widehat{Y}$ means that material, geometry, and time inputs are mapped to a predicted response. The scope is directional response, not complete field-scale geological simulation.

\KeyPoints
\begin{itemize}
  \item The task is defined by a source point and a probe point.
  \item Distance, unit direction, and azimuth are explicit geometric quantities.
  \item The prediction maps material, geometry, and time inputs to selected outputs.
\end{itemize}

\Transition
To train and compare the models, the work uses a controlled synthetic dataset.
\end{slidenote}

\begin{slidenote}{COMSOL-Based Synthetic Dataset}
\Purpose
Explain where the reference data come from and why they are synthetic.

\WhatToSay
The training data is synthetic, generated in COMSOL Multiphysics with a coupled thermoelastic model. The setup is a 2D plane-strain domain, one by one metre, with an initial temperature of 293 kelvin and a localized source at 1500 kelvin. We export temperature, the two displacement components, stress, strain, and velocity as reference fields for all surrogates.

\KeyPoints
\begin{itemize}
  \item The dataset is COMSOL-derived and synthetic.
  \item The domain is a simplified 2D plane-strain setting.
  \item Exported fields provide reference data for surrogate comparison.
\end{itemize}

\Transition
The next slide shows the mesh and numerical domain used for these exports.
\end{slidenote}

\begin{slidenote}{COMSOL 2D Model and Mesh}
\Purpose
Show the numerical domain and explain why mesh structure matters.

\WhatToSay
The geometry is a 2D section of rock with an internal heated source. The mesh is refined near the source, where temperature, displacement, and stress gradients are strongest. This mesh defines the spatial structure from which all fields are exported.

\KeyPoints
\begin{itemize}
  \item The simulation is a simplified two-dimensional domain.
  \item Mesh refinement is concentrated around the thermal source.
  \item The mesh provides the spatial structure for exported fields.
\end{itemize}

\Transition
Now I will show the main simulated thermoelastic fields.
\end{slidenote}

\begin{slidenote}{COMSOL Simulation Results}
\Purpose
Connect the synthetic thermal source with temperature, displacement, and stress outputs.

\WhatToSay
Three representative fields. Local heating creates a compact temperature disturbance near the source. Thermal expansion then produces a displacement field around the heated region. And the stress concentrates near the source. Together they confirm the coupled thermoelastic behaviour we expect, and they are the ground truth the models learn from.

\KeyPoints
\begin{itemize}
  \item Local heating produces a localized temperature disturbance.
  \item Thermal expansion creates displacement.
  \item Stress concentration confirms coupled thermoelastic behaviour in the synthetic setup.
\end{itemize}

\Transition
After simulation, the exported fields must be represented in a model-ready format.
\end{slidenote}

\begin{slidenote}{Data Representation and Preprocessing}
\Purpose
Explain how COMSOL exports are converted into inputs for different model families.

\WhatToSay
COMSOL exports are aligned into a common node-based representation. Dynamic fields are assembled into a spatio-temporal tensor of time steps, points, and field channels; static material features are stored separately. Every channel is normalized before training. The same processed dataset is reused by the coordinate-based, grid-based, and graph-based models -- that's what makes the comparison fair.

\KeyPoints
\begin{itemize}
  \item CSV exports are aligned by nodes and time steps.
  \item Dynamic fields form a spatio-temporal tensor.
  \item Normalization makes model training and comparison more stable.
\end{itemize}

\Transition
Next, I will describe the comparative prototype framework.
\end{slidenote}

\begin{slidenote}{Comparative Prototype Framework}
\Purpose
Explain that all model routes are evaluated under one shared workflow.

\WhatToSay
All four models run through one shared workflow: the same frontend, backend orchestrator, material catalogue, and a single input--output contract. Same source--probe geometry, same material presets, same normalized response -- temperature, displacement, direction, travel time, and diagnostics. This is what lets us compare model families rather than implementations.

\KeyPoints
\begin{itemize}
  \item All routes use comparable inputs.
  \item Outputs are normalized into the same response structure.
  \item The framework supports fair model comparison.
\end{itemize}

\Transition
Before discussing individual models, I will define the shared comparison contract.
\end{slidenote}

\begin{slidenote}{Shared Model Comparison Contract}
\Purpose
Clarify the inputs, outputs, and interpretation shared by all surrogate routes.

\WhatToSay
Each route is treated as a surrogate for the same prediction task, not as a full solver. The input is material parameters, geometry, time, and field features. The output is temperature, displacement, travel time, and diagnostics. The comparison basis is identical presets, the same response format, and one analytical travel-time sanity check.

\KeyPoints
\begin{itemize}
  \item The comparison uses one shared prediction contract.
  \item Inputs and outputs are kept consistent across models.
  \item The routes are surrogate models, not equivalent physical solvers.
\end{itemize}

\Transition
Before discussing the individual models, I will briefly describe the implementation stack.
\end{slidenote}

\begin{slidenote}{Implementation Stack}
\Purpose
Explain the software stack that connects the web interface with the model routes.

\WhatToSay
The prototype is a local Docker Compose workspace. A web frontend for scenario setup; a FastAPI backend that validates requests and normalizes responses; and separate Python services for the PINN, FNO, MeshGraphNet, and Transformer routes. The backend does not replace the models -- it orchestrates them and returns one comparable response format.

\KeyPoints
\begin{itemize}
  \item Docker Compose runs the local demo stack.
  \item FastAPI provides validation, routing, and response normalization.
  \item Python model services expose PINN, FNO, MeshGraphNet, and Transformer routes.
\end{itemize}

\Transition
With the implementation stack defined, I will start the model discussion with the physics-informed baseline.
\end{slidenote}

\begin{slidenote}{PINN: Physics-Informed Baseline}
\Purpose
Introduce the PINN route and explain why it is the main physical baseline.

\WhatToSay
PINN is the explicitly physics-informed route. It trains on supervised data plus thermoelastic residuals and predicts the temperature and displacement fields. It serves as the physics-consistency baseline against which the data-driven routes are judged. It is a neural surrogate whose loss encourages physical consistency -- not an exact solver.

\KeyPoints
\begin{itemize}
  \item PINN combines supervised data with physical residual terms.
  \item It predicts temperature and displacement fields.
  \item It is the strongest physics-informed baseline in the comparison.
\end{itemize}

\Transition
The next slide shows the PINN input contract and architecture.
\end{slidenote}

\begin{slidenote}{PINN Input and Architecture}
\Purpose
Explain what the PINN receives, what it predicts, and why the architecture is differentiable.

\WhatToSay
The network is a multilayer perceptron: an input layer of physical features, five hidden blocks of width 192 with tanh activation, and an output layer for the predicted fields. Inputs are geometry, time, and material parameters; outputs are temperature and displacement. Velocity, stress, and strain are recovered through automatic differentiation during residual evaluation.

\KeyPoints
\begin{itemize}
  \item Inputs are geometry, time, and material properties.
  \item The implemented network is $9\rightarrow192$, five hidden $192\rightarrow192$ blocks, and $192\rightarrow3$ output.
  \item Differentiability allows residuals and velocity consistency to be evaluated from derivatives.
\end{itemize}

\Transition
With the architecture defined, I will explain the composite loss.
\end{slidenote}

\begin{slidenote}{PINN Composite Loss and Residuals}
\Purpose
Show how data fitting and physical consistency are combined in the PINN route.

\WhatToSay
The objective is a weighted sum of four terms: a supervised field loss, a velocity-consistency term from time derivatives, an elastic-wave residual, and a coupled thermal residual. All residuals are computed from automatic derivatives. The current limitation is that dedicated boundary- and initial-condition losses are not yet enforced.

\KeyPoints
\begin{itemize}
  \item The loss combines supervised, velocity, elastic, and thermal terms.
  \item Residuals are computed using automatic differentiation.
  \item Boundary and initial losses are a current limitation.
\end{itemize}

\Transition
The second model family is the Fourier Neural Operator.
\end{slidenote}

\begin{slidenote}{Fourier Neural Operator}
\Purpose
Introduce FNO as the structured-grid neural operator route.

\WhatToSay
The FNO learns an operator between fields rather than point values. The input is a multi-channel field on a grid, and the output is the predicted thermoelastic field -- three channels: temperature and two displacement components. Because it learns in function space, it captures global spatial patterns efficiently.

\KeyPoints
\begin{itemize}
  \item FNO learns a mapping between spatial fields.
  \item The practical setup uses two-dimensional tensors.
  \item The output channels represent temperature and in-plane displacement.
\end{itemize}

\Transition
The next slide shows the main FNO processing pipeline.
\end{slidenote}

\begin{slidenote}{FNO Pipeline}
\Purpose
Summarize the processing stages of the FNO route.

\WhatToSay
The field is first lifted into latent features, passed through a stack of spectral blocks, then projected back to the output field. Inside each spectral block: a Fourier transform into frequency space, multiplication of retained modes by trainable weights, an inverse transform, and a parallel local convolution. So it combines global spectral interactions with local spatial corrections.

\KeyPoints
\begin{itemize}
  \item FNO uses lifting, spectral processing, and projection.
  \item Fourier operations capture global spatial structure.
  \item Local convolution helps represent short-range details.
\end{itemize}

\Transition
Now I will explain the spectral block more directly.
\end{slidenote}

\begin{slidenote}{FNO Spectral Block}
\Purpose
Explain the Fourier-space operation and retained modes at an intuitive level.

\WhatToSay
Each layer update sums a spectral branch and a local convolution branch, then applies an activation. The spectral branch transforms the field, keeps only the low-frequency modes, multiplies them by a trainable weight tensor, and transforms back. Low modes carry the smooth global structure; the local convolution restores short-range detail. The number of retained modes is a key hyperparameter.

\KeyPoints
\begin{itemize}
  \item The spectral branch works in Fourier space.
  \item Retained modes control spectral resolution and cost.
  \item The block is data-driven, not a PDE residual calculation.
\end{itemize}

\Transition
The next slide gives the training objective and the current caveat for FNO.
\end{slidenote}

\begin{slidenote}{Training and Evaluation of FNO}
\Purpose
State how FNO is trained and make the current calibration limitation transparent.

\WhatToSay
The pipeline is dataset, normalization, model, prediction, loss, evaluation. The loss combines mean-squared error with a relative L2 term. Normalization statistics come from the training set and are reused everywhere. Evaluation uses field comparisons, channel-wise errors, and physical sanity checks. The current caveat: the FNO output scale still needs recalibration before it is physically reliable.

\KeyPoints
\begin{itemize}
  \item FNO training uses supervised MSE and relative error terms.
  \item Channel normalization is reused consistently.
  \item The current output scale needs recalibration.
\end{itemize}

\Transition
The third model family is MeshGraphNet.
\end{slidenote}

\begin{slidenote}{MeshGraphNet: Graph Formulation}
\Purpose
Introduce MeshGraphNet as the mesh-aware surrogate route.

\WhatToSay
MeshGraphNet represents the COMSOL mesh directly as a graph -- nodes are mesh points, edges connect neighbours. Each node carries its coordinates, current physical state, material parameters, and scenario condition. Each edge carries the relative geometry between two nodes. The model predicts the next physical state for every node.

\KeyPoints
\begin{itemize}
  \item The mesh is represented as nodes and edges.
  \item Node features carry state, material, and scenario information.
  \item Edge features carry local geometric relations.
\end{itemize}

\Transition
The next slide shows the input and output structure in more detail.
\end{slidenote}

\begin{slidenote}{MeshGraphNet: Inputs and Outputs}
\Purpose
Clarify what is stored at nodes and edges and what the graph model predicts.

\WhatToSay
On the left, the two inputs: node feature vectors with position, state, material, and condition; and edge attributes with relative position and distance. The mesh becomes a graph $G$ of vertices and edges. The output is the predicted next-step field vector for every node.

\KeyPoints
\begin{itemize}
  \item Nodes contain geometry, material data, and physical state.
  \item Edges describe local spatial relationships.
  \item The model predicts next physical fields over mesh nodes.
\end{itemize}

\Transition
Now I will summarize the graph neural network architecture.
\end{slidenote}

\begin{slidenote}{MeshGraphNet: Encoder--Processor--Decoder}
\Purpose
Describe the MeshGraphNet architecture without excessive implementation detail.

\WhatToSay
Three stages. Encoders map node and edge features into a latent space of dimension 128. Then ten message-passing blocks propagate information between neighbouring nodes -- this is the processor. Finally a node-wise decoder maps the latent state to the predicted output fields.

\KeyPoints
\begin{itemize}
  \item Encoder: embeds raw node and edge features.
  \item Processor: performs repeated message passing.
  \item Decoder: converts latent node states to physical outputs.
\end{itemize}

\Transition
The next slide explains the message-passing block itself.
\end{slidenote}

\begin{slidenote}{MeshGraphNet: Message Passing Block}
\Purpose
Connect graph message passing with local physical interaction learning.

\WhatToSay
Each block has two updates. First the edge update: every edge is refined from its two endpoint nodes and its own features, with a residual connection and layer norm. Then messages are aggregated per destination node by scatter-add, and the node update refines each node from its state and the aggregated messages. The residual connections keep repeated message passing stable.

\KeyPoints
\begin{itemize}
  \item Edge updates exchange information between neighbouring nodes.
  \item Aggregation collects incoming messages at each node.
  \item Physical behaviour is learned from data and must be validated.
\end{itemize}

\Transition
The fourth model family is the Transformer-style surrogate.
\end{slidenote}

\begin{slidenote}{Transformer: Inputs and Outputs}
\Purpose
Introduce the attention-based route and its token representation.

\WhatToSay
The transformer is a neural operator on a point set. Input tokens, on the left, live in R-thirteen: coordinates, initial state, and material parameters. Query tokens, on the right, are 2D coordinates where we want a prediction. At each query the model returns R-three: temperature and two displacement components. Input and query sizes are independent, so the operator is discretisation-invariant.

\KeyPoints
\begin{itemize}
  \item Input tokens contain state, geometry, and material information.
  \item Query tokens specify prediction locations.
  \item Outputs are temperature and in-plane displacement at query points.
\end{itemize}

\Transition
Next, I will show how the Transformer architecture processes those tokens.
\end{slidenote}

\begin{slidenote}{Transformer: Architecture}
\Purpose
Summarize the encoder--decoder attention mechanism for query prediction.

\WhatToSay
Two pipelines. Encoder: input tokens projected into the embedding space, then self-attention layers -- each token attends to every other. Decoder: queries projected and passed through cross-attention layers where keys and values come from the encoder output. A linear head maps decoded queries back to $T$, $u$, $v$.

\KeyPoints
\begin{itemize}
  \item Encoder self-attention processes known input tokens.
  \item Decoder cross-attention connects query points with encoded inputs.
  \item The output head predicts $(T,u,v)$ at requested points.
\end{itemize}

\Transition
The next slide defines the training loss for this route.
\end{slidenote}

\begin{slidenote}{Transformer: Training Loss}
\Purpose
Explain the supervised objective and its limited physical regularisation.

\WhatToSay
The loss is a weighted sum of three terms with weights one, zero-point-two-five, and zero-point-zero-five. First, supervised MSE against COMSOL ground truth. Second, velocity consistency: the time derivative of predicted displacement, obtained by autodiff, matches the exported velocity. Third, a thermal residual that enforces the diffusion side of the heat equation. Trained with AdamW, token subsampling for bounded attention cost.

\KeyPoints
\begin{itemize}
  \item The objective combines supervised, velocity, and thermal terms.
  \item Token subsampling controls attention cost.
  \item Physics constraints are weaker than in PINN.
\end{itemize}

\Transition
Now I will move from model descriptions to calculation status and results.
\end{slidenote}

\begin{slidenote}{Calculation Status}
\Purpose
Summarize the evaluation setup and state caveats before showing plots.

\WhatToSay
The comparison covers 40 scenarios per material, four materials, four models -- 640 active responses, zero fallback, all in 2D plane strain. One caveat to state up front: the FNO output scale is still miscalibrated, so we treat it as a diagnostic branch rather than a ranked predictor.

\KeyPoints
\begin{itemize}
  \item The comparison covers $640$ model responses.
  \item Plotted outputs come from active models.
  \item FNO is included but currently miscalibrated.
\end{itemize}

\Transition
The first result slide compares predicted travel times.
\end{slidenote}

\begin{slidenote}{Travel-Time Predictions}
\Purpose
Explain what the travel-time plot shows and how to interpret the model differences.

\WhatToSay
Mean predicted travel time per route: PINN around 0.11 milliseconds, Transformer 0.14, FNO 1.3, MeshGraphNet 2.5. The key observation is that the differences between model families are much larger than the difference between basalt and sandstone -- in this prototype the model choice dominates the material effect.

\KeyPoints
\begin{itemize}
  \item The plot compares travel-time estimates across models and materials.
  \item Model-family differences dominate the material difference here.
  \item The result motivates careful validation.
\end{itemize}

\Transition
To evaluate accuracy, the next slide compares predictions with an analytical reference.
\end{slidenote}

\begin{slidenote}{Accuracy Against Analytical Reference}
\Purpose
Present the main quantitative result and its cautious interpretation.

\WhatToSay
We sanity-check against the analytical body-wave travel time, distance over P-velocity. PINN has a median error of two percent and stays inside the ten-percent band for every scenario. Transformer is about thirty percent, FNO thirty-five, MeshGraphNet over a hundred. So PINN is the only route that is quantitatively trustworthy on travel time today.

\KeyPoints
\begin{itemize}
  \item The analytical reference is a sanity check, not complete validation.
  \item PINN shows the best agreement with the reference.
  \item Other routes need calibration or stronger time conditioning.
\end{itemize}

\Transition
Accuracy is one side of the comparison; inference latency is also important.
\end{slidenote}

\begin{slidenote}{Inference Latency and Time Behavior}
\Purpose
Compare practical response time and explain the fixed-horizon issue.

\WhatToSay
On speed the order reverses: Transformer is fastest at about 1.3 milliseconds, then MeshGraphNet, FNO, and PINN at 2.6. All four are well under ten milliseconds. One more point: PINN is the only route that varies smoothly with the requested observation time; the others currently behave as fixed-horizon predictors.

\KeyPoints
\begin{itemize}
  \item All routes are fast enough for interactive prototype comparison.
  \item Transformer has the lowest median latency.
  \item PINN currently has the most meaningful time-conditioned behaviour.
\end{itemize}

\Transition
Before concluding, I will state the main threats to validity.
\end{slidenote}

\begin{slidenote}{Threats to Validity}
\Purpose
Show scientific caution and define the limits of the results.

\WhatToSay
Four honest limitations. The training data is a pilot COMSOL configuration, not a calibrated geological database. The setup is simplified 2D with effective parameters. The FNO normalization is not yet reliable. And the analytical reference is isotropic -- a sanity check, not a full thermoelastic validation.

\KeyPoints
\begin{itemize}
  \item The dataset is limited and synthetic.
  \item The physical formulation is simplified.
  \item The analytical reference is useful but limited.
\end{itemize}

\Transition
After the limitations, I will briefly show the web application prototype.
\end{slidenote}

\begin{slidenote}{Web Application Overview}
\Purpose
Introduce the web interface as the user-facing part of the prototype.

\WhatToSay
The whole comparison is exposed through a web application. This is the landing view: a research-prototype workspace for directional prediction in a fixed 2D domain.

\KeyPoints
\begin{itemize}
  \item The web interface is the user-facing part of the prototype.
  \item It supports model comparison under one workflow.
  \item The prototype scope is stated directly in the interface.
\end{itemize}

\Transition
Next, I will show the full prediction workspace.
\end{slidenote}

\begin{slidenote}{Directional Prediction Workspace}
\Purpose
Show the full interface layout used after a prediction is returned.

\WhatToSay
The user picks a material and a model route, sets the source and probe coordinates inside the unit square, and chooses the observation time. The form maps one-to-one onto the shared prediction contract that the backend validates.

\KeyPoints
\begin{itemize}
  \item The workspace collects response panels and runtime information in one screen.
  \item The output is normalized into a common format for all routes.
  \item This supports comparison rather than manual inspection of raw service payloads.
\end{itemize}

\Transition
Now I will zoom into the source--probe geometry visualization.
\end{slidenote}

\begin{slidenote}{Source--Probe Visualization}
\Purpose
Explain how the interface visualizes the directional prediction query.

\WhatToSay
The workspace draws the geometry live: source and probe points, the propagation arrow with distance and azimuth, and a unit-direction key. The geometry shown here is computed by the same backend logic used in the prediction, so the picture and the numbers always agree.

\KeyPoints
\begin{itemize}
  \item The source and probe define the directional query.
  \item Distance, azimuth, and unit direction are shown visually.
  \item The visualization connects model output with the geometry of the task.
\end{itemize}

\Transition
The next slide shows the response dashboard in more detail.
\end{slidenote}

\begin{slidenote}{Prediction Response Dashboard}
\Purpose
Show the main numerical response panels returned by the selected model.

\WhatToSay
After a run, the response is shown as cards: the thermal response, the displacement components and magnitude, the directional summary, and the travel time -- the same normalized format for every model route.

\KeyPoints
\begin{itemize}
  \item The dashboard shows thermal, mechanical, directional, and temporal outputs.
  \item Model metadata makes the selected route and latency visible.
  \item The same response structure is used across routes.
\end{itemize}

\Transition
The final interface slide shows runtime metadata and diagnostics.
\end{slidenote}

\begin{slidenote}{Runtime Diagnostics and Metadata}
\Purpose
Show how the prototype reports service status, metadata, and limitations.

\WhatToSay
At the bottom, the diagnostics strip reports the run state, the model version, inference time, request identifier, and any warnings. This keeps every prediction traceable and makes clear when a route is a real model output versus a fallback.

\KeyPoints
\begin{itemize}
  \item Diagnostics expose model identity, latency, and request metadata.
  \item Warnings and notes make runtime assumptions visible.
  \item The interface reinforces the prototype scope before the conclusion.
\end{itemize}

\Transition
With the prototype shown, I can summarize the conclusions.
\end{slidenote}

\begin{slidenote}{Conclusions}
\Purpose
State the main scientific outcomes of the diploma work.

\WhatToSay
Four takeaways. We built a physically grounded comparative framework for this prediction task. PINN is the strongest physics-informed baseline, with the best agreement against the analytical reference. Transformer is the fastest route, and MeshGraphNet is a useful mesh-aware candidate. FNO, for now, is a diagnostic branch until its output normalization is fixed.

\KeyPoints
\begin{itemize}
  \item The main contribution is a comparative AI-assisted framework.
  \item PINN performs best against the travel-time sanity check.
  \item Other routes remain useful candidates with specific limitations.
\end{itemize}

\Transition
Finally, I will outline the most important next steps.
\end{slidenote}

\begin{slidenote}{Future Work}
\Purpose
Show how the prototype can be developed after the diploma.

\WhatToSay
Next: retrain and recalibrate the 2D FNO; add time-conditioned rollout to MeshGraphNet; extend the material catalogue with more complete thermoelastic and validation data; and compare against additional analytical or numerical baselines. To be precise about scope -- the prototype estimates directional characteristics within the simplified assumptions of the study.

\KeyPoints
\begin{itemize}
  \item Recalibrate FNO.
  \item Improve time-conditioned behaviour in MeshGraphNet.
  \item Extend material data and validation baselines.
\end{itemize}

\Transition
This completes the main presentation.
\end{slidenote}

\begin{slidenote}{Thank You}
\Purpose
Close the presentation and invite questions.

\WhatToSay
Thank you for your attention -- we are happy to take questions.

\KeyPoints
\begin{itemize}
  \item Close politely.
  \item Signal readiness to discuss assumptions and limitations.
  \item Keep answers consistent with the research-prototype framing.
\end{itemize}

\Transition
Questions from the committee.
\end{slidenote}

\begin{slidenote}{Backup: Mechanical Rock Parameters}
\Purpose
Use this backup only if the committee asks about mechanical material values.

\WhatToSay
This backup table gives the extended mechanical material catalogue. It includes density, Young's modulus, Poisson's ratio, shear modulus, bulk modulus, and representative P-wave and S-wave velocities for ten rock types. I would not read every value. The main point is that the prototype uses physically meaningful effective parameters, while still treating them as midpoint values under isotropic assumptions.

\KeyPoints
\begin{itemize}
  \item Mechanical parameters control inertia, stiffness, and wave speeds.
  \item Values are effective and comparative, not site-calibrated.
  \item The table supports the material assumptions used in the main flow.
\end{itemize}

\Transition
If needed, the next backup slide gives thermal and porosity parameters.
\end{slidenote}

\begin{slidenote}{Backup: Thermal and Porosity Parameters}
\Purpose
Use this backup if asked how thermal properties and porosity enter the catalogue.

\WhatToSay
This table gives thermal expansion, thermal conductivity, heat capacity, volumetric heat capacity, thermoelastic coupling, and porosity information. These quantities control how a temperature disturbance spreads and how temperature change contributes to thermal strain and stress. Dashes indicate values that were not available in the source catalogue. They are left as missing values rather than filled by unsupported assumptions.

\KeyPoints
\begin{itemize}
  \item Thermal properties control heat spreading and thermoelastic strain.
  \item Porosity is included as an effective descriptive property.
  \item Missing values are shown explicitly.
\end{itemize}

\Transition
If the question concerns wave velocities, use the next backup slide.
\end{slidenote}

\begin{slidenote}{Backup: Elastic Wave Velocities}
\Purpose
Explain the P-wave and S-wave formulas if the committee asks about wave physics.

\WhatToSay
This slide gives the standard formulas for compressional and shear wave velocities. P-waves involve particle motion mainly parallel to the propagation direction and include volume change. S-waves involve transverse motion and depend on shear rigidity. The formulas show that stiffness increases wave velocity, while density increases inertia. These relations also explain why the analytical reference $t=r/V_p$ is a useful first sanity check for travel time.

\KeyPoints
\begin{itemize}
  \item $V_P$ depends on bulk modulus, shear modulus, and density.
  \item $V_S$ depends on shear modulus and density.
  \item These formulas justify the simple travel-time reference.
\end{itemize}

\Transition
If the question concerns heat-spreading time scale, use the next backup slide.
\end{slidenote}

\begin{slidenote}{Backup: Thermal Diffusivity}
\Purpose
Explain the compact thermal parameter used to interpret heat spreading.

\WhatToSay
Thermal diffusivity is defined as $\alpha_T=k/(\rho C_p)$. It has units of square metres per second and describes how quickly a temperature disturbance spreads through a material. Larger thermal diffusivity means faster temperature smoothing. This quantity is useful because it combines conductivity, density, and heat capacity into one interpretable parameter for the thermal part of the coupled response.

\KeyPoints
\begin{itemize}
  \item Thermal diffusivity describes heat-spreading rate.
  \item It combines conductivity, density, and heat capacity.
  \item It helps interpret the time scale of the thermal response.
\end{itemize}

\Transition
The final backup slides document image credits and references.
\end{slidenote}

\begin{slidenote}{Image Credits}
\Purpose
Document visual sources if the committee asks about images.

\WhatToSay
This slide lists the sources for the external educational diagrams and rock textures used in the presentation. It helps separate illustrative visual material from the thesis results. The COMSOL field images and model-result plots are part of the project workflow, while the external diagrams are used only to support physical explanation.

\KeyPoints
\begin{itemize}
  \item External images are credited separately.
  \item Illustrative diagrams are not model outputs.
  \item Project-generated figures are distinct from source images.
\end{itemize}

\Transition
The reference slide gives the scientific and technical sources behind the work.
\end{slidenote}

\begin{slidenote}{References}
\Purpose
Point to the scientific basis of the thesis if asked.

\WhatToSay
The references include rock physics, thermoelasticity, geodynamics, surface geophysics, anisotropy, neural operators, and Transformer-related sources. In the defence, I would use this slide only in the question period. It supports the theoretical background, material-property assumptions, and machine-learning model choices used throughout the thesis.

\KeyPoints
\begin{itemize}
  \item References cover physics, geology, and neural surrogate modelling.
  \item They support the theory and material assumptions.
  \item The slide is mainly for Q\&A support.
\end{itemize}

\Transition
End of backup material.
\end{slidenote}

\newpage
\section*{\color{AcademicNavy}Possible questions from the committee}
\addcontentsline{toc}{section}{Possible questions from the committee}

\textbf{1. Why did you choose directional prediction instead of full-field simulation?}

Because the thesis focuses on a narrower and more defensible research task. Full-field thermoelastic simulation is more computationally and physically complex. Directional prediction allows us to compare source--probe responses, travel time, and model behaviour under controlled assumptions.

\textbf{2. Is the prototype validated on real field data?}

No. The current work uses COMSOL-derived synthetic data and analytical sanity checks. The results should be interpreted within simplified modelling assumptions. Field data would be an important next step, but it is outside the current scope.

\textbf{3. Why is PINN called a physics-informed baseline?}

PINN is called a physics-informed baseline because its training objective includes residual terms connected with thermoelastic relations. This means that physical structure is included in the learning objective, but the result is still interpreted as a baseline model inside the simplified prototype.

\textbf{4. Why does FNO perform poorly in the current results?}

The current two-dimensional FNO route has a calibration and output-scale issue. Therefore, it is kept in the comparison as a diagnostic branch. Its architecture is still relevant, but this implementation needs retraining and normalization correction.

\textbf{5. What is the role of the analytical reference $t=r/V_p$?}

It is a simple travel-time sanity check based on distance and P-wave velocity. It is not a complete thermoelastic validation, but it helps identify whether predicted travel times are physically plausible in a basic isotropic approximation.

\textbf{6. Why is Transformer fast but less physically reliable?}

The Transformer route uses attention-based query prediction and can return estimates quickly, but most of its behaviour is learned from data. It has weaker explicit physical constraints than PINN, so its predictions require careful validation.

\textbf{7. What is the current limitation of MeshGraphNet?}

MeshGraphNet preserves mesh topology and is suitable for mesh-aware surrogate modelling, but in the current prototype it behaves mainly as a fixed-horizon predictor. It needs stronger time-conditioned rollout behaviour for more reliable time-dependent prediction.

\textbf{8. What are the main limitations of the work?}

The main limitations are the simplified two-dimensional setting, effective isotropic material parameters, limited synthetic data, fixed-horizon behaviour in some surrogate routes, the FNO normalization issue, and the absence of direct calibration against real geological measurements.

\textbf{9. What is the main contribution of the diploma work?}

The main contribution is a physically grounded comparative framework for AI-Guided Prediction of Thermoelastic Wave Propagation in Geological Media. It combines thermoelastic theory, geological material parameters, a controlled dataset, and several neural surrogate routes under one consistent comparison workflow.

\newpage
\section*{\color{AcademicNavy}Short version of the speech}
\addcontentsline{toc}{section}{Short version of the speech}

Good morning. The topic of our diploma work is \textit{AI-Guided Prediction of Thermoelastic Wave Propagation in Geological Media}. The work presents a research prototype for AI-guided directional prediction of thermoelastic wave propagation in simplified geological media. It is not a field-calibrated simulator. The goal is to estimate selected directional characteristics and compare several neural surrogate approaches under one shared contract.

The physical motivation is that temperature changes in rocks can create thermal expansion, stress, displacement, and wave-like mechanical response. This is relevant for geothermal systems, underground engineering, and geophysical interpretation. A full coupled numerical simulation can be expensive when many materials and source--probe directions must be compared, so the thesis formulates a narrower directional prediction task.

The model uses effective geological material parameters such as density, elastic moduli, thermal conductivity, heat capacity, thermal expansion, and P-wave and S-wave velocities. The formulation is two-dimensional, assumes plane strain, small deformation, and linear isotropic thermoelasticity, and does not explicitly resolve pores, fluids, fractures, plasticity, or full anisotropy.

Directional prediction is defined by a source point and a probe point. Their difference gives the direction, distance, and azimuth. The prototype estimates temperature, displacement, travel time, and directional response for that query. The data are based on a controlled COMSOL-derived two-dimensional thermoelastic setup, exported and normalized for several model families.

The comparison includes a Physics-Informed Neural Network, Fourier Neural Operator, MeshGraphNet, and Transformer-style surrogate. PINN is used as the physics-informed baseline because its objective includes thermoelastic residual terms. FNO, MeshGraphNet, and Transformer-style models are treated as surrogate routes for comparison, not as physically equivalent numerical solvers.

The calculations use forty scenarios per material, four materials, and four model routes, giving six hundred forty model responses. In the travel-time comparison, model-family differences are larger than the material difference in the current prototype. Against the simple analytical reference $t=r/V_p$, PINN shows the best agreement and remains within the ten percent band in the balanced sweep. Transformer-style prediction is the fastest route, while MeshGraphNet and FNO need additional development for time behaviour and calibration.

The main limitations are the limited synthetic dataset, the simplified two-dimensional assumptions, the current FNO normalization issue, fixed-horizon behaviour in some routes, and the fact that the analytical reference is only a sanity check. The conclusion is that the work provides a physically grounded comparative framework for AI-guided thermoelastic wave prediction within simplified assumptions. Future work should recalibrate FNO, improve time-conditioned MeshGraphNet behaviour, extend material data, and compare against additional analytical or numerical baselines.

\end{document}
